\documentclass[12pt]{article}

\usepackage{fouriernc}
\usepackage{amsmath}

\usepackage{tikz}

\newcommand{\contradiction}{%
\begin{tikzpicture}[rotate=45,x=0.5ex,y=0.5ex]
\draw[line width=.2ex] (0,2) -- (3,2) (0,1) -- (3,1) (1,3) -- (1,0) (2,3) -- (2,0);
\end{tikzpicture}
}

\begin{document}

\section*{Problem-183}
\subsection*{Problem Statement}
The first term of an arithmetic progression is $a$, the 4th term is $b$. Find the second and the third terms.
\subsection*{Solution}
The $n$-th term of the arithmetic progression has the form $a+(n-1)d$ where $a$ is the first term ($n=1$) and $d$ is the difference. From the information given in the statement, we have:
\[
a + (4-1)d = b
\]
Thus $d = \frac{b-a}{3}$. So, the second term is $a+(2-1)d = a+\frac{b-a}{3} = \frac{2a+b}{3}$. The third term is $d$ more than the second term, that is, $\frac{a+2b}{3}$.

\section*{Problem-198}
\subsection*{Problem Statement}
The first term of a geometric progression is $a$ and the third term is $b$. Find the second term.
\subsection*{Solution}
Say the second term is $x$. Then from the definition of a geometric progression, $\frac{x}{a} = \frac{b}{x}$. Therefore, $x^2 = ab$. We shall consider three separate possibilities.
\begin{enumerate}
\item When $ab < 0$, there is no real-valued $x$. In other words, there is no such real-valued geometric progression. For example, even if the ratio were negative like -1, the geometric progression would alternate the signs, so that first and third terms would still have the same sign.
\item When $ab = 0$, we have $x = 0$.
\item When $ab > 0$, there are two possibilities for $x$: $\sqrt{ab}$ and $-\sqrt{ab}$. 
\end{enumerate}

\section*{Problem-204}
\subsection*{Problem Statement}
Is it possible that numbers $\frac{1}{2}$, $\frac{1}{3}$, and $\frac{1}{5}$ are (not necessarily adjacent) terms of the same arithmetic progression?

\subsection*{Solution}
If the numbers come from the same arithmetic progression, then we can consider them as sorted. Since, in the question, the numbers are already in descending order, we shall go with that order. So, if the three numbers indeed come from the same arithmetic progression, the difference will turn out to be negative. We may consider $\frac{1}{2}$ as the first term of the sub-progression. We thus have:
\begin{equation*}
\begin{aligned}
\frac{1}{2} + (m-1)\ d &= \frac{1}{3}\\
\frac{1}{2} + (n-1)\ d &= \frac{1}{5}
\end{aligned}
\end{equation*}
Where $m \neq 1$, $m < n$ and $d\ (\neq 0)$ is the difference. We can rearrange and get the below two equations:
\begin{equation}
		(m-1)\ d = -\frac{1}{6}	
\end{equation}
\begin{equation}
		(n-1)\ d = -\frac{3}{10}	
\end{equation}
Dividing $(1)$ by $(2)$, we have:
\begin{equation}
	\frac{m-1}{n-1} = \frac{5}{9}
\end{equation}
Choosing $(m, n) = (6, 10)$ agrees with $(3)$. Using $m=6$ in $(1)$ gives $d = -\frac{1}{30}$. So, the numbers could be a part of the same arithmetic progression because, $\frac{1}{2} + 5\ \left( -\frac{1}{30} \right) = \frac{1}{3}$ and $\frac{1}{2} + 9 \ \left( -\frac{1}{30} \right) = \frac{1}{5}$. Hence, starting from $\frac{1}{2}$, the 5-th term is $\frac{1}{3}$ and the 9-th term is $\frac{1}{5}$. Actually, there are other valid choices like $(m, n, d) = \left( 11, 19, -\frac{1}{60} \right)$.

\section*{Problem-205}
\subsection*{Problem Statement}
Is it possible that the numbers 2, 3, and 5 are (not necessarily adjacent) terms of a geometric progression?

\subsection*{Solution}
If the numbers belong to the same geometric progression, they must be sorted. Since the numbers are in increasing order in the question, we shall go with that order. So, for the sub-progression that starts at 2, we could consider 2 as the first term. We thus have the following:
\setcounter{equation}{0}
\begin{equation*}
	\begin{aligned}
		2 \cdot q^m &= 3\\
		q^{mn} &= \frac{3^n}{2^n}
	\end{aligned}
\end{equation*}
And,
\begin{equation*}
	\begin{aligned}
		2 \cdot q^n &= 5\\
		     q^{mn} &= \frac{5^m}{2^m}
	\end{aligned}
\end{equation*}
Here $0 < m < n$ and $q$ is the ratio. We can then write:
\begin{equation*}
	\begin{aligned}
		\frac{3^n}{2^n} &= \frac{5^m}{2^m}\\
		3^n &= 2^{n-m} \cdot 5^m
	\end{aligned}
\end{equation*}
On the left side we have an odd number and on the right side we have an even number. We have reached a contradiction. \contradiction\\

Therefore, the numbers cannot be part of the same geometric progression. Note, we would reach a contradiction even if we considered the numbers in reverse order.

\section*{Problem-210}
\subsection*{Problem Statement}
Is it possible that the second term of a geometric progression is less than its first term and also less than its third term?

\subsection*{Solution}
Let's assume that the first three terms of the geometric progression are $a$, $aq$, and $aq^2$. When $a = 0$ or $q = 0$, the answer is negative. So, we shall now consider $a \neq 0$ and $q \neq 0$ case. We shall consider two separate possibilities.
\begin{enumerate}
\item When $a < 0$. If $aq < a$, then $q > 1$. If it is also true that $aq < aq^2$, then $q > q^2$. Both $q > 1$ and $q > q^2$ cannot be true at the same time.
\item When $a > 0$. If $aq < a$, we have $q < 1$. If it is also true that $aq < aq^2$, then $q < q^2$. Both $q < 1$ and $q < q^2$ are possible for say $q = -1$.
\end{enumerate}

So, $1,\ -1,\ 1, \ldots$ is a possible geometric progression.

\section*{Problem-212}
\subsection*{Problem Statement}
Is it possible that an infinite arithmetic progression contains exactly two integer terms?

\subsection*{Solution}
It is not possible. We shall use \textbf{proof by contradiction}. Therefore, we shall assume that exactly two terms of an infinite arithmetic progression are integers and that would lead us to a contradiction, completing the proof.

Say $a+md = x$ and $a+nd = y$ are the only two integral terms where $a$ is the first term of the progression, $0 < m < n$, and $d$ is the difference. We notice that $(n-m)d = y-x$ is an integer. Now consider the term $y + (n-m)d = z$. $z$ must come after $y$, so $z$ is another term distinct from $x$ and $y$. By our assumption, $z$ must not be an integer. However, $y$ and $(n-m)d$ are both integers, so should be their sum, i.e., $z$. We have our contradiction. \contradiction 

\end{document}